Maia-3 (Chessformer) is the strongest open square-token model for predicting human chess moves. We
reproduce its evaluation protocol exactly and its reported accuracies to within a third of a point,
benchmark all nine released checkpoints, and audit the paper and code: we find that its public engine
runs without the move history it was trained with (costing 2.2 points of accuracy), that its history
encoding and Elo conditioning waste most of the input, that its published models ignore the clock,
and that its headline ablation omits the baseline needed to credit its dynamic attention bias.
We then present \textbf{Cloverb}, a square-token transformer conditioned on exactly what a player
has at the board: their own rating, the time control, the time left on their clock, and the last
64 plies encoded as moves. A simple prior, initialising static attention biases with chess movement
patterns, makes Cloverb learn 5.5 points faster than the Maia-3 architecture at equal data, and in
these 512k-position diagnostics a dynamic bias term in the style of Chessformer's adds nothing.
With a C data pipeline, a fused AVX-512 attention kernel and bf16 training on Intel AMX, all
experiments, including the final model distilled from Maia-3 through amortised teacher pools, ran
on 4-core CPUs without a GPU.
On the ALLIE benchmark under Maia's protocol the final 9.2M-parameter Cloverb reaches 55.3\% top-1 and 82.4\% top-3 move matching (Maia-3-5M: 55.3\%, Maia-3-79M: 56.8\% on the same positions).
